% Figure: isentropic area-Mach relation A/A* versus M for gamma=1.4.
% Non-monotonic: large at low M, minimum 1 at M=1, large again at high M.
% The same area ratio admits one subsonic and one supersonic solution.
\begin{figure}[htbp]
\centering
\begin{tikzpicture}
\begin{axis}[
  width=12cm, height=8cm,
  xlabel={Mach number $M$}, ylabel={area ratio $A/A^{*}$},
  xmin=0, xmax=4, ymin=0, ymax=11,
  axis lines=left,
  grid=major, grid style={enggray!20},
  tick label style={font=\scriptsize},
  label style={font=\small},
  legend style={font=\scriptsize, at={(0.97,0.97)}, anchor=north east},
]
% A/A* = (1/M)*[ (2/(g+1))(1 + (g-1)/2 M^2) ]^((g+1)/(2(g-1)))
% g=1.4: exponent = 2.4/0.8 = 3.0; factor 2/2.4 = 0.833333
% A/A* = (1/M)*(0.833333*(1+0.2*M^2))^3
\addplot[engblue, very thick, domain=0.08:4, samples=200] {
  (1/x)*(0.833333*(1+0.2*x^2))^3 };
% M=1 marker at A/A*=1
\addplot[engred, only marks, mark=*, mark options={fill=engred}] coordinates {(1,1)};
\node[engred, font=\scriptsize, anchor=south west] at (axis cs:1.05,1.1) {$M=1$, $A/A^{*}=1$};
% example: A/A*=2 admits two solutions; horizontal line
\draw[enggray, dashed] (axis cs:0,2) -- (axis cs:4,2);
\node[enggray, font=\scriptsize, anchor=south] at (axis cs:0.3,2.05) {$A/A^{*}=2$};
% the two intersections of A/A*=2: M approx 0.31 (subsonic), 2.20 (supersonic)
\addplot[engblack, only marks, mark=square, mark options={draw=engblack}] coordinates {(0.31,2) (2.20,2)};
\node[engblack, font=\scriptsize, anchor=west] at (axis cs:0.33,2.5) {subsonic};
\node[engblack, font=\scriptsize, anchor=west] at (axis cs:2.25,2.5) {supersonic};
\end{axis}
\end{tikzpicture}
\caption[Isentropic area-Mach relation]{The area-Mach relation
$A/A^{*}(M)$ for $\gamma=1.4$. The curve is non-monotonic: it
diverges as $M\to 0$ (a slow subsonic flow needs a large area), takes
its minimum value of $1$ at $M=1$ (the throat), and rises again as
$M\to\infty$ (a fast supersonic flow needs a large area because the
density has dropped). A given area ratio above $1$ therefore admits
two isentropic solutions, one subsonic and one supersonic; for
$A/A^{*}=2$ these are $M\approx 0.31$ and $M\approx 2.20$. Which
branch a nozzle realises is set by the back pressure. The plotted
curve is the exact relation; the two square markers are the roots at
$A/A^{*}=2$.}
\label{fig:vol03ch13:area-mach}
\end{figure}
